% Storage-engine overhead on the ingest path — \input{} into evaluation_dynamic.tex,
% after dynamic_v2.tex (which reports the ingest throughput this section attributes).
% Assumes wiredtiger_basics.tex has already introduced connections, sessions, cursors,
% the cache, skip lists, update chains, and reconciliation.

\section{Isolating the Persistence Machinery Tax}
\label{sec:wt-ingest-overhead}
\newflag[new: Results from profiling the per-edge overhead when running an ingest-only workload]

\project{} uses \wt{} as its storage engine because \wt{} provides crash recovery, transactions, and a mature concurrency model.
When a workload does not need persistence, \project{} disables write-ahead logging (\texttt{log=(enabled=false)}), checkpoints (\texttt{checkpoint=(wait=0, log\_size=0)}), and raises all eviction thresholds to 100\% so that no eviction work occurs.
These settings eliminate disk I/O on the write path.

What CPU overhead does \wt{} still impose on the ingest path once we turn persistence off?
Answering this question establishes the performance penalty \project{} pays for using a general-purpose transactional storage engine rather than a purpose-built in-memory index.

\subsection{Persistence Components Are Negligible}
\label{sec:wt-overhead-persistence}

We answer the question with a combination of source-code analysis and profiling.
We profiled the same ingest-only experiment we report in \autoref{sec:dynamic:ingest-update-tput}: graph500-24 ingestion at one and 72 writers with \texttt{perf}, resolving each sample to its \wt{} source function through DWARF.
Profiling this experiment shows that \project{} spends 13.8\,$\mu$s of writer CPU time per inserted edge at one writer, 3.8$\times$ that of Teseo on the same workload.
We first show where that time does \emph{not} go, then where it does.

Only the checkpoint code path and the eviction code path call reconciliation.
The cursor insert path appends an update to the page's skip list and returns; no reconciliation call occurs on this path.
Block allocation is reachable only through the block-write path invoked at the end of reconciliation.
Only eviction triggers page splits, and the split itself is purely in-memory: it allocates a new page, migrates insert-list entries, and updates parent pointers, with no disk writes and no block allocation.

With WAL and checkpoints disabled, these persistence components together account for 46\,ns per inserted edge, under 0.5\% of writer CPU time.

\subsection{Where the Time Goes}
\label{sec:wt-overhead-costs}

The costs that dominate the ingest path are intrinsic to \wt{}'s B-tree access path and the bookkeeping that accompanies every page access.
\autoref{tab:wt-overhead} shows the per-edge breakdown at one and 72 writers on graph500-24, with WAL and checkpoints disabled.

\begin{table}[h]
  \centering
  \begin{tabularx}{\columnwidth}{@{}Xrrrr@{}}
    \toprule
    Component & w=1 (ns) & \% & w=72 (ns) & Growth \\
    \midrule
    B-tree descent and insertion           & 5,421 & 39 & 6,664 & $\times$1.2 \\
    Key comparison during descent          & 2,180 & 16 & 2,603 & $\times$1.2 \\
    Cache accounting and eviction protocol & 1,047 & 8 & 18,879 & $\times$18 \\
    Cursor, transaction, config parsing    & 1,819 & 13 & 13,017 & $\times$7 \\
    Memory allocation and page faults      & 1,570 & 11 & 653 & $\times$0.4 \\
    Key serialization (packing)            & 185 & 1 & 264 & $\times$1.4 \\
    Persistence components                 & 46 & $<$1 & 62 & $\times$1.3 \\
    Other                                  & 1,554 & 11 & 3,520 & $\times$2.3 \\
    \midrule
    Writer on-CPU total & \textbf{13,822} & & \textbf{45,662} & $\times$\textbf{3.3} \\
    \bottomrule
  \end{tabularx}
  \caption{On-CPU time per inserted edge for \project{} on graph500-24 with WAL and
  checkpoints disabled.
  Persistence components include reconciliation, block management, WAL, and page splits.
  The ``other'' row includes generation counters, dirty-page tracking, libc string
  operations, kernel overhead, and unresolved samples.}
  \label{tab:wt-overhead}
\end{table}

\paragraph{Tree operations.}
B-tree descent and in-memory skip-list insertion account for 39\% of single-writer CPU time; byte-string key comparison during descent adds another 16\%, making tree operations the largest single cost at 55\%.
Any storage engine that maintains a sorted B-tree over variable-length keys would pay a comparable cost for descent and comparison; this is not specific to \wt{}'s persistence design.
Key \emph{serialization}, the packing step that produces the byte-comparable encoding, accounts for only 1.3\%.
The format conversion itself is cheap; searching with the result is not.

\paragraph{Cache accounting and eviction protocol.}
This row is architecturally distinct from the other costs.
\wt{} manages a \emph{bounded cache}: data lives on disk, and it can evict pages from memory and read them back later.
The entire eviction protocol exists to support this design.
Every piece of per-access bookkeeping in this row traces back to it: cache byte counters track memory usage so eviction knows when to trigger; threshold checks compare those counters against eviction targets on every page access; hazard pointers protect pages from being evicted while a thread is reading them; and the LRU read generation scores pages for eviction priority.

A pure in-memory index eliminates all this bookkeeping.
Teseo, for example, has no page cache, no eviction, and no hazard pointers, and pays zero
in this category.
Yet this protocol runs on every \wt{} page access even when we configure the eviction
thresholds never to trigger.

At one writer, the protocol costs 1,047\,ns per edge (8\%).
At 72 writers, it grows to 18,879\,ns (41\%), an 18$\times$ increase, because every insert
both writes and reads the same connection-wide cache byte counters.

\wt{} creates the eviction server thread unconditionally at connection startup;
no configuration can disable it.
It runs a wake-check-sleep loop with an auto-scaling interval, but with thresholds at
100\% it finds nothing to do and returns to sleep.
The cost measured above comes from the per-access protocol that application threads
execute, not the background thread.

\paragraph{API and transaction bookkeeping.}
Cursor setup and teardown, transaction management (snapshot walks over active sessions,
transaction-ID allocation), and \wt{}'s configuration-string parser together account for
13\% at one writer.
Under contention these grow sharply: the cursor API alone reaches 6,346\,ns at 72 writers
($\times$23 its single-writer cost), and transaction bookkeeping grows to 5,169\,ns
($\times$13).

\paragraph{Memory allocation.}
At one writer, memory allocation and kernel page-fault handling account for 11\% of CPU time.
Much of the page-fault cost is kernel work rather than \wt{} code.
At 72 writers this row drops to 653\,ns because the faults amortize over a shorter per-edge wall-clock time.

\subsection{Implications}
\label{sec:wt-overhead-implications}

The ingest-path overhead that \project{} pays for using \wt{} with persistence disabled
falls into three categories, summarized in \autoref{tab:wt-overhead-categories}.

\begin{table}[h]
  \centering
  \begin{tabularx}{\columnwidth}{@{}XrrX@{}}
    \toprule
    Category & w=1 & w=72 & Why it exists \\
    \midrule
    Persistence operations (reconciliation, block mgmt, WAL, page splits)
      & $<$1\% & $<$1\% & Persisting data to disk. Disabled and negligible. \\
    \addlinespace
    Persistence architecture (cache accounting, eviction protocol, hazard pointers)
      & 8\% & 41\% & Bounded-cache design. Runs on every page access even when eviction
        thresholds prevent any eviction from occurring. \\
    \addlinespace
    B-tree access path (descent, key comparison, transactions, cursor API, key packing)
      & $\sim$81\% & $\sim$51\% & Operating a transactional sorted index. Any storage
        engine with a B-tree, MVCC, and a cursor API would pay a comparable cost. \\
    \bottomrule
  \end{tabularx}
  \caption{Three categories of ingest-path CPU overhead when \wt{}'s persistence is
  disabled.
  Percentages are shares of single-writer (w=1) and 72-writer (w=72) on-CPU time, computed
  from \autoref{tab:wt-overhead}; the balance in each column is the ``other'' row.}
  \label{tab:wt-overhead-categories}
\end{table}

Persistence \emph{operations} are negligible once we turn WAL and checkpoints off.
The B-tree access path is a cost any transactional sorted index would impose.
The middle category is the one specific to \wt{}'s persistence \emph{architecture}: the
bounded-cache design requires per-access bookkeeping (cache byte counters, threshold
checks, hazard pointers, LRU scoring) that a pure in-memory index does not need.
At one writer this tax is modest (8\%), but under contention it grows to 41\% and becomes
the largest single component of per-edge cost, ahead of B-tree descent and of the cursor
and transaction layer, because all threads contend on the same connection-wide counters.

\project{} accepts this overhead because it can re-enable \wt{}'s crash-recovery
guarantees without changing the data representation.
The bounded cache and its eviction protocol are the mechanism that makes this possible:
they keep the in-memory data organized in pages that \wt{} can reconcile and write to
disk on demand.
